% Presentación abreviada y precisa del contenido de un documento, sin agregar interpretación o crítica; no debe exceder una página.
\section*{Resumen}

Muchas plataformas de empleo buscan comparando las palabras de la consulta con las de la oferta, y fallan cuando ambas describen el mismo cargo con términos distintos. Este trabajo evalúa en qué medida un pipeline de generación aumentada por recuperación (RAG) mejora la recuperación de ofertas frente a la búsqueda por palabras clave, qué componentes explican esa mejora y con qué costo en latencia. Siguiendo Design Science Research, el artefacto es el pipeline de JobAI, una plataforma preexistente de empleo remoto, que se mejoró en dos iteraciones y se evaluó con un protocolo repetible: un corpus congelado de 1.068 ofertas reales, 100 consultas en inglés y en español, BM25 con y sin los mismos filtros, y juicios de relevancia de un modelo de lenguaje validado contra juicios humanos ($\kappa = 0{,}79$).

La configuración completa obtuvo un nDCG@10 de 0,82, frente a 0,56 de BM25 y 0,60 de BM25 con filtros. El reordenamiento fue lo que más aportó (+0,10) y lo que más costó, unos 0,8~s por consulta; la búsqueda híbrida no aportó. Con los filtros, la proporción de los cinco primeros resultados que cumplían las restricciones pasó de 0,41 a 0,95, y en español la configuración completa alcanzó 0,76, frente a 0,33 de BM25. La elección del motor vectorial solo importó con filtros selectivos a gran escala. La fidelidad de las respuestas (0,90) no cambió entre iteraciones, aunque la primera ignoraba las restricciones: las métricas de RAGAS sobre la respuesta final no detectaron esa falla, que sí mostró la medida de cada componente.

\textbf{Palabras clave}: generación aumentada por recuperación, búsqueda de empleo, recuperación de información, búsqueda híbrida, reordenamiento, bases de datos vectoriales, evaluación, Design Science Research.

\newpage
\section*{Abstract}

Many job platforms search by matching the words of the query against those of the posting, and fail when both describe the same position in different terms. This work evaluates to what extent a retrieval-augmented generation (RAG) pipeline improves job retrieval over keyword search, which components explain that improvement, and at what latency cost. Following Design Science Research, the artifact is the pipeline of JobAI, a pre-existing remote job search platform, which was improved in two iterations and evaluated with a repeatable protocol: a frozen corpus of 1,068 real job postings, 100 queries in English and Spanish, BM25 with and without the same filters, and relevance judgments from a language model validated against human judgments ($\kappa = 0.79$).

The full configuration reached an nDCG@10 of 0.82, against 0.56 for BM25 and 0.60 for BM25 with filters. Reranking was the component that contributed the most (+0.10) and cost the most, about 0.8~s per query; hybrid search did not help. With the filters, the share of the top five results that met the constraints rose from 0.41 to 0.95, and in Spanish the full configuration reached 0.76, against 0.33 for BM25. The choice of vector engine mattered only with selective filters at large scale. Answer faithfulness (0.90) did not change between iterations, even though the first one ignored the constraints: RAGAS metrics on the final answer did not detect that failure, which the per-component measurement did reveal.

\textbf{Keywords}: retrieval-augmented generation, job search, information retrieval, hybrid search, reranking, vector databases, evaluation, Design Science Research.
